\documentclass[12pt,a4paper]{article}

\usepackage[T1]{fontenc}
\usepackage[utf8]{inputenc}
\usepackage{amsmath,amssymb,amsthm}
\usepackage{mathtools}
\usepackage{physics}
\usepackage{booktabs}
\usepackage{array}
\usepackage{colortbl}
\usepackage{xcolor}
\usepackage{hyperref}
\usepackage{geometry}
\usepackage{setspace}
\usepackage{fancyhdr}
\usepackage{tcolorbox}
\usepackage{enumitem}
\usepackage[numbers]{natbib}
\tcbuselibrary{theorems,skins,breakable}

\geometry{margin=2.5cm, top=3cm, bottom=3cm}
\onehalfspacing

% ── Colours ──────────────────────────────────────────────────────────────────
\definecolor{tuogold}{gray}{0.0}
\definecolor{tuoteal}{gray}{0.0}
\definecolor{tuored}{gray}{0.0}
\definecolor{tuogreen}{gray}{0.0}
\definecolor{tuopurple}{gray}{0.0}
\definecolor{tuoblue}{gray}{0.0}
\definecolor{lightgold}{gray}{0.95}
\definecolor{lightteal}{gray}{0.95}
\definecolor{lightred}{gray}{0.95}
\definecolor{lightgreen}{gray}{0.95}
\definecolor{lightpurple}{gray}{0.95}
\definecolor{lightblue}{gray}{0.95}

% ── tcolorbox environments ────────────────────────────────────────────────────
\tcbset{enhanced, breakable, sharp corners, boxrule=0.8pt, left=8pt, right=8pt}

\newtcolorbox{resultbox}[1]{
  colback=lightgold, colframe=black,
  title={\textbf{Result: #1}},
  fonttitle=\small\bfseries, coltitle=black}

\newtcolorbox{openprobbox}[1]{
  colback=lightblue, colframe=black,
  title={\textbf{Open Problem #1}},
  fonttitle=\small\bfseries, coltitle=black}

\newtcolorbox{verdictbox}{
  colback=lightgreen, colframe=black,
  title={\textbf{Verdict}},
  fonttitle=\small\bfseries}

\newtcolorbox{warningbox}{
  colback=lightred, colframe=black,
  title={\textbf{Null result}},
  fonttitle=\small\bfseries}

\newtcolorbox{keybox}{
  colback=lightpurple, colframe=black,
  title={\textbf{Key Insight}},
  fonttitle=\small\bfseries}

% ── Theorem environments ──────────────────────────────────────────────────────
\newtheorem{theorem}{Theorem}[section]
\newtheorem{corollary}[theorem]{Corollary}
\newtheorem{lemma}[theorem]{Lemma}
\newtheorem{proposition}[theorem]{Proposition}
\theoremstyle{definition}
\newtheorem{definition}[theorem]{Definition}
\newtheorem{remark}[theorem]{Remark}

% ── Macros ────────────────────────────────────────────────────────────────────
\newcommand{\tPl}{t_{\mathrm{Pl}}}
\newcommand{\lPl}{\ell_{\mathrm{Pl}}}
\newcommand{\TPl}{T_{\mathrm{Pl}}}
\newcommand{\EPl}{E_{\mathrm{Pl}}}
\newcommand{\TTUO}{T_{\mathrm{TUO}}}
\newcommand{\gstar}{g_{*}}
\newcommand{\Ngen}{N_{\mathrm{gen}}}
\newcommand{\VPl}{V_{\mathrm{Pl}}}
\newcommand{\ZSC}{\mathcal{Z}_{\mathrm{ZS}}}
\newcommand{\Zcorr}{Z_{\mathrm{corr}}}
\newcommand{\Sact}{S_{\mathrm{E}}}
\newcommand{\Ffree}{F_{\mathrm{free}}}
\newcommand{\Hmfp}{\lambda_{\mathrm{mfp}}}
\newcommand{\alphas}{\alpha_{s}}
\newcommand{\etaobs}{\eta_{\mathrm{obs}}}
\newcommand{\eps}{\varepsilon}
\newcommand{\dBmL}{\delta(B{-}L)}

% ── Header / footer ───────────────────────────────────────────────────────────
\pagestyle{fancy}
\fancyhf{}
\renewcommand{\headrulewidth}{0.4pt}
\fancyhead[L]{\small TUO — Open Problems Analysis}
\fancyhead[R]{\small Matshaba (2026)}
\fancyfoot[C]{\thepage}

\hypersetup{
  colorlinks=true, linkcolor=black,
  citecolor=black, urlcolor=black,
  pdftitle={TUO Open Problems: Numerical Analysis},
  pdfauthor={Romeo Matshaba}}

% ─────────────────────────────────────────────────────────────────────────────
\begin{document}

\begin{titlepage}
\centering
\vspace*{2cm}
{\LARGE\bfseries Theory of Universal Origins\\[0.5em]
\Large Open Problems: Numerical Analysis\\[0.3em]
\large OP5 · OP1 · OP2 · OP4}
\vspace{1.5em}

\rule{\textwidth}{0.4pt}
\vspace{0.5em}

{\large Romeo Matshaba}\\[0.3em]
{\normalsize Department of Physics, University of South Africa (UNISA)\\
Pretoria, South Africa}\\[0.5em]
{\normalsize March 2026}\\[0.3em]
{\small DOI: \href{https://doi.org/10.5281/zenodo.18827993}{10.5281/zenodo.18827993}}

\vspace{1em}\rule{\textwidth}{0.4pt}\vspace{1.5em}

\begin{abstract}
We carry out detailed numerical analysis of four open problems in the
Theory of Universal Origins (TUO).
The central new result concerns OP5 (generation selection): we prove that the
constrained path integral weight $Z_{\mathrm{corr}}(g_*)$, dissected into
three layers---Euclidean action, thermal entropy, and Gaussian charge
fluctuations---does not select $\Ngen = 3$.
The free energy $F = E - T_{\mathrm{TUO}} S$ is strictly decreasing in $\Ngen$
with no minimum; the $Z_{\mathrm{corr}}$ correction is of order $O(1)$, corresponding
to $\sim 10^{-93}$ per cell, and is completely negligible against the thermal
entropy. The resolution is a two-sided pinch:\ the Kobayashi--Maskawa theorem
provides the hard lower bound $\Ngen \geq 3$ (CP violation required for
$\eta \neq 0$), while the LEP measurement $N_\nu = 2.984 \pm 0.008$ provides
the upper bound $\Ngen \leq 3$. Consequently OP5 and OP1 (baryon asymmetry)
are the same open problem.
For OP4 (junction time), we improve the upper bound from $94\,\tPl$
(isothermal) to $152\,\tPl$ using the FRW cooling law $T(t)\propto t^{-1/2}$.
For OP2 (CMB), we establish the suppression budget: the $2.4\%$ cell-level
fluctuation $\delta T/T = 1/(4\sqrt{\gstar})$ must be reduced by
$\sim 10^4$ via Sachs--Wolfe freezing and acoustic propagation.
\end{abstract}
\end{titlepage}

\tableofcontents
\newpage

% ─────────────────────────────────────────────────────────────────────────────
\section{Introduction}
% ─────────────────────────────────────────────────────────────────────────────

The Theory of Universal Origins derives the initial conditions of the Hot Big
Bang from two axioms:\ a flat Minkowski background and the global zero-sum
constraint $\mathrm{Tr}[\hat\rho(t)\,\hat Q_k] = 0$ for all conserved charges
$\hat Q_k$. From these alone, fourteen quantitative results follow without free
parameters, including $w = 1/3$, $H(t_{\mathrm{Pl}}) = 1/(2\tPl)$, $\Omega = 1$,
and $T_{\mathrm{TUO}} = (15/\pi^2)^{1/4} T_{\mathrm{Pl}}$.

Six open problems remain. Dark energy (OP3) is the hardest and is not
addressed here. This paper examines the four tractable problems:

\begin{itemize}[leftmargin=2em]
\item \textbf{OP5}: Does the constrained path integral select $\Ngen = 3$?
\item \textbf{OP1}: What mechanism produces $\eta = 6.12\times10^{-10}$?
\item \textbf{OP2}: What sets the CMB power spectrum amplitude?
\item \textbf{OP4}: What is the precise TUO--FRW junction time?
\end{itemize}

We treat OP5 first because, as we show, it is the gateway to all others.
All numerical results use the cube volume convention $\VPl = \lPl^3$,
CODATA\,2018 constants, and PDG\,2024 particle data.
The complete Python implementation is in \texttt{tuo\_open\_problems\_results.py}
(same repository, same DOI).

% ─────────────────────────────────────────────────────────────────────────────
\section{SM Content and the g* Counting Convention}
\label{sec:gstar}
% ─────────────────────────────────────────────────────────────────────────────

The effective number of relativistic degrees of freedom $\gstar$ is fundamental
to TUO. \emph{All numerical results use the standard cosmological convention
that includes both particles and antiparticles.}

\begin{table}[h]
\centering
\caption{SM degree-of-freedom counting per generation
(standard convention: particles + antiparticles).}
\label{tab:gstar_counting}
\begin{tabular}{lrrr}
\toprule
Sector & Modes & SB weight & $\gstar$ contribution \\
\midrule
Photon $\gamma$ & 2 & 1 & 2.000 \\
Gluons $g$ & 16 & 1 & 16.000 \\
$W^+, W^-$ & $3+3$ & 1 & 6.000 \\
$Z^0$ & 3 & 1 & 3.000 \\
Higgs $H$ & 1 & 1 & 1.000 \\
\midrule
Boson subtotal & 28 & — & 28.000 \\
\midrule
Quarks $u, \bar{u}$ per gen & $6+6$ & $7/8$ & 10.500 \\
Quarks $d, \bar{d}$ per gen & $6+6$ & $7/8$ & 10.500 \\
Leptons $e, \bar{e}$ per gen & $2+2$ & $7/8$ & 3.500 \\
Leptons $\nu_L, \bar\nu_L$ per gen & $1+1$ & $7/8$ & 1.750 \\
\midrule
Fermion per-gen subtotal & 30 & $7/8$ & 26.250 \\
\midrule
Total ($\Ngen = 3$) & $28 + 90 = 118$ & — & $28 + 3\times26.25 = \mathbf{106.75}$ \\
\bottomrule
\end{tabular}
\end{table}

The general formula is $\gstar(\Ngen) = 28 + 26.25\,\Ngen$, giving
$\gstar(1) = 54.25$, $\gstar(2) = 80.50$, $\gstar(3) = 106.75$.
The zero-sum constraint requires $B - L = 0$ and $Q = 0$ at species level
(without spin multiplicity), which holds for all $\Ngen$ independently.

% ─────────────────────────────────────────────────────────────────────────────
\section{OP5: The Z\textsubscript{corr} Calculation}
\label{sec:op5}
% ─────────────────────────────────────────────────────────────────────────────

\subsection{Setup}

The constrained path integral for the emergence probability is:
\begin{equation}
  \mathcal{P}(\Ngen) = \mathcal{N}^{-1} \int_{\ZSC} \mathcal{D}\hat\rho\;
  e^{-S_E[\hat\rho]},
\end{equation}
where $\ZSC$ is the intersection of zero-sum hyperplanes. At saddle point,
$S_E = \gstar(\Ngen)/2$ (the Heisenberg minimum for all modes). The
Gaussian fluctuation integral around the saddle gives:
\begin{equation}
  \ln \mathcal{P}(\Ngen) = -\Sact(\Ngen) + \ln Z_1(\Ngen) + \ln\Zcorr(\Ngen),
\end{equation}
where $Z_1$ is the single-cell canonical partition function and
$\Zcorr$ encodes the global constraint.

The thermal quantities at $E = \EPl/2$, $T = \TTUO$ are:
\begin{equation}
  x \equiv \frac{\EPl}{2k_B\TTUO} = \frac{\EPl}{2k_B(15/\pi^2)^{1/4}\TPl}
  = 0.4503.
\end{equation}
Since $x < 1$, all modes are in the classical regime. The occupation
numbers are:
\begin{equation}
  n_B = \frac{1}{e^x-1} = 1.758, \qquad
  n_F = \frac{1}{e^x+1} = 0.389.
\end{equation}

\subsection{Layer 1: Euclidean Action}

\begin{equation}
  \Sact(\Ngen) = \frac{\gstar(\Ngen)}{2}
  = 14 + 13.125\,\Ngen \quad [\hbar].
\end{equation}
This is strictly increasing in $\Ngen$, so the action alone strongly prefers
$\Ngen = 1$ (minimum action = $27.125\,\hbar$).

\subsection{Layer 2: Thermal Entropy}

Each bosonic mode contributes $S_B = 1.806\,k_B$ and each fermionic mode
$S_F = 0.668\,k_B$ to the single-cell entropy. The single-cell
log-partition function is:
\begin{equation}
  \ln Z_1(\Ngen) = N_B\ln z_B + N_F(\Ngen)\ln z_F
  = 28\ln(2.758) + 30\,\Ngen\ln(1.637).
\end{equation}
The increment per generation:
\begin{equation}
  \frac{\partial \ln Z_1}{\partial \Ngen} = 30\ln z_F = 14.79 > 0.
\end{equation}
The free energy per cell:
\begin{equation}
  \Ffree(\Ngen) = \Sact - \frac{\TTUO}{\TPl}\,S_{\rm th}
  = \Sact - 1.110\,S_{\rm th},
\end{equation}
decreases by $9.14\,\EPl$ per added generation.
\emph{Entropy outweighs action for all $\Ngen$.}

\begin{table}[h]
\centering
\caption{Three-layer decomposition of $\ln\mathcal{P}(\Ngen)$ per cell.}
\label{tab:layers}
\begin{tabular}{rrrrrrrc}
\toprule
$\Ngen$ & $\gstar$ & $\Sact$ & $T\cdot S_{\rm th}$ & $\Ffree$ &
$\ln Z_1$ & $\ln\Zcorr$ & CP phases \\
\midrule
1 & 54.25 & 27.13 & 78.42 & $-$51.29 & 43.20 & $-$1.68 & 0 \\
2 & 80.50 & 40.25 & 100.68 & $-$60.43 & 57.99 & $-$2.41 & 0 \\
\rowcolor{lightgold}
3 & 106.75 & 53.38 & 122.95 & $-$69.57 & 72.79 & $-$2.85 & \textbf{1} \\
4 & 133.00 & 66.50 & 145.21 & $-$78.71 & 87.58 & $-$3.17 & 3 \\
5 & 159.25 & 79.63 & 167.48 & $-$87.85 & 102.38 & $-$3.43 & 6 \\
6 & 185.50 & 92.75 & 189.74 & $-$96.99 & 117.17 & $-$3.64 & 10 \\
\bottomrule
\end{tabular}
\end{table}

\subsection{Layer 3: Z\textsubscript{corr} Gaussian Correction}

The global zero-sum constraint introduces a Gaussian suppression
$\Zcorr \propto (\det\chi)^{-1/2}$, where $\chi$ is the $3\times3$
charge susceptibility matrix for $(Q_{\rm em}, B, L)$:
\begin{equation}
  \ln\Zcorr(\Ngen) = -\tfrac{1}{2}\ln\det\chi(\Ngen).
\end{equation}
The susceptibilities are:
\begin{align}
  \chi_Q &= 6\chi_B^{(W)} + \Ngen\bigl(8\cdot(4/9) + 12\cdot(1/9) + 4\bigr)\chi_F
  \approx 34 + 1.27\,\Ngen, \\
  \chi_B &= \Ngen\cdot\frac{24}{9}\chi_F \approx 1.27\,\Ngen, \\
  \chi_L &= 6\Ngen\,\chi_F \approx 1.43\,\Ngen,
\end{align}
giving $\ln\Zcorr$ of order $O(1)$.

For the $N = 8.8\times10^{92}$ cells in the observable universe, the
per-cell correction is:
\begin{equation}
  \frac{\ln\Zcorr(\Ngen)}{N_{\rm cells}} \sim 10^{-93}.
\end{equation}
The total entropy gain per generation across all cells is
$\Delta\ln Z = N_{\rm cells}\times 14.79 \sim 10^{94}$,
which is $10^{94}$ times larger. The $\Zcorr$ correction is
completely negligible.

\begin{warningbox}
The constrained path integral $\mathcal{P}(\Ngen)$ has no maximum.
The free energy $\Ffree(\Ngen)$ is strictly decreasing.
$Z_{\rm corr}$ does \emph{not} select $\Ngen = 3$.
\end{warningbox}

\subsection{Resolution: Two-Sided Pinch}

The question of why $\Ngen = 3$ is not answered by the path integral
weight alone. It requires two independent inputs:

\subsubsection{Lower bound: Kobayashi--Maskawa theorem}

The $\Ngen\times\Ngen$ CKM matrix contains
$\frac{(\Ngen-1)(\Ngen-2)}{2}$ physical CP-violating phases:
\begin{equation}
  \delta_{\rm CP} \neq 0 \quad\Longleftrightarrow\quad \Ngen \geq 3.
\end{equation}
For $\Ngen \leq 2$: the CKM matrix is real, $\eps_{\rm CP} = 0$,
leptogenesis gives $\eta = 0$, and no matter-dominated universe exists.
The constraint $\eta_{\rm obs} \neq 0$ is therefore equivalent to $\Ngen \geq 3$.

\subsubsection{Upper bound: LEP measurement}

The decay width of the $Z^0$ boson into invisible channels (light neutrinos)
gives:
\begin{equation}
  N_\nu = 2.984 \pm 0.008 \quad\Rightarrow\quad \Ngen \leq 3.
\end{equation}

\subsubsection{Pinch}

\begin{equation}
  \{\Ngen \geq 3\} \cap \{\Ngen \leq 3\} = \{3\}.
\end{equation}

\begin{verdictbox}
$\Ngen = 3$ exactly, from the intersection of the Kobayashi--Maskawa
lower bound and the LEP upper bound.
TUO provides the thermal bath in which leptogenesis operates.
The KM theorem then enforces $\Ngen \geq 3$ as a
\emph{precondition for a matter-dominated universe to exist}.
\end{verdictbox}

\begin{keybox}
\textbf{OP5 and OP1 are the same open problem.}
Selecting $\Ngen$ and explaining $\eta$ require the same
mechanism---CP violation in the lepton sector---applied to the
same TUO thermal bath.
Closing OP1 (full leptogenesis Boltzmann calculation at $T_{\rm TUO}$)
simultaneously closes OP5.
\end{keybox}

% ─────────────────────────────────────────────────────────────────────────────
\section{OP1: Baryon Asymmetry}
\label{sec:op1}
% ─────────────────────────────────────────────────────────────────────────────

\subsection{Initial State}

At emergence, the TUO configuration contains:
\begin{equation}
  B = L = 2\Ngen = 6 \quad \text{(at species level, no antiparticles)},
\end{equation}
\begin{equation}
  (B-L)_{\rm init} = 0 \quad \text{exactly (Axiom II).}
\end{equation}
The photon number per Planck cell at $T = \TTUO$ is:
\begin{equation}
  N_\gamma = \frac{2\zeta(3)}{\pi^2}\left(\frac{k_B\TTUO}{\hbar c}\right)^3 \VPl
  = 0.3334.
\end{equation}
The naive ratio $\eta_{\rm naive} = B/N_\gamma = 18$ is $\sim3\times10^{10}$
times too large and has the wrong sign for standard leptogenesis---it represents
the \emph{initial} baryon content before sphaleron washout.

\subsection{Sphaleron Washout}

Electroweak sphalerons are active for $T > T_{\rm EW} \approx 130\,\mathrm{GeV}$.
They conserve $B-L$ but violate $B+L$.
For exact $(B-L) = 0$, sphalerons drive $B+L \to 0$ and hence $\eta \to 0$
at tree level. This is not a problem; it is the correct starting point.

\subsection{Required CP Asymmetry}

The SM sphaleron conversion gives $B = \tfrac{28}{79}(B-L)$.
To produce $\eta_{\rm obs} = 6.12\times10^{-10}$, we need:
\begin{equation}
  \dBmL = \frac{79}{28}\,\eta_{\rm obs}\,N_\gamma
  = 1.73\times10^{-9}.
\end{equation}
In the simplified strong-washout Boltzmann treatment with efficiency $\kappa = 0.05$:
\begin{equation}
  \eps_{\rm CP} \approx \frac{\eta_{\rm obs}\,\gstar}{\kappa}
  \approx 1.3\times10^{-6}.
\end{equation}
The Davidson--Ibarra bound on the right-handed neutrino CP asymmetry is:
\begin{equation}
  \eps_{\rm max}(M_N) = \frac{3}{16\pi}
  \frac{M_N\,\Delta m_\nu^2}{v_{\rm EW}^2}
  \geq \eps_{\rm CP},
\end{equation}
which requires:
\begin{equation}
  M_N \geq M_N^{\rm min}
  = \frac{16\pi\,v_{\rm EW}^2\,\eps_{\rm CP}}{3\,\Delta m_\nu^2}
  \approx 3\times10^{12}\,\mathrm{GeV}.
\end{equation}

\subsection{The KM Link}

\begin{theorem}[Kobayashi--Maskawa CP Violation Threshold]
\label{thm:km}
The CKM matrix for $\Ngen$ generations has
$(\Ngen-1)(\Ngen-2)/2$ physical CP-violating phases.
For $\Ngen = 1$ or $2$: the matrix is real, all CP phases vanish,
the Davidson--Ibarra asymmetry $\eps_{\rm CP} = 0$, and
leptogenesis produces $\eta = 0$.
For $\Ngen \geq 3$: CP violation is kinematically allowed.
\end{theorem}

The implication is immediate: any universe with $\Ngen \leq 2$ and the
TUO initial state $(B-L) = 0$ is matter-symmetric at all times.
The observation $\eta_{\rm obs} > 0$ therefore implies $\Ngen \geq 3$
as a strict logical consequence.

\begin{resultbox}{OP1 -- OP5 equivalence}
$\eta \neq 0 \;\Leftrightarrow\; \eps_{\rm CP} \neq 0 \;\Leftrightarrow\;
\Ngen \geq 3$.
The open problem ``explain $\eta$'' contains the open problem ``select $\Ngen$''
as a prerequisite. They are closed by the same calculation:
a full leptogenesis Boltzmann system at $T = \TTUO$ with the CP-asymmetric
RH neutrino decay rate as input.
\end{resultbox}

% ─────────────────────────────────────────────────────────────────────────────
\section{OP2: CMB Power Spectrum}
\label{sec:op2}
% ─────────────────────────────────────────────────────────────────────────────

\subsection{Cell-Level Fluctuation}

In a single Planck cell containing $\gstar$ modes, the energy fluctuation
by the central limit theorem is:
\begin{equation}
  \frac{\delta E}{E} = \frac{1}{\sqrt{\gstar}} = \frac{1}{\sqrt{106.75}}
  = 0.0968 \quad (9.7\%).
\end{equation}
Since $E \propto T^4$:
\begin{equation}
  \frac{\delta T}{T}\bigg|_{\rm cell} = \frac{1}{4}\frac{\delta E}{E}
  = \frac{1}{4\sqrt{\gstar}} = 0.0242 \quad (2.4\%).
\end{equation}
The observed amplitude is $\delta T/T \sim 10^{-5}$: a suppression of
$\sim 2400$, or approximately $11$ binary e-folds.

\subsection{Sachs--Wolfe Transfer}

The gravitational potential at horizon crossing is:
\begin{equation}
  \Phi \approx \frac{\delta E}{E}\left(\frac{\tPl}{t}\right)^2,
\end{equation}
since the quantum correction $\Delta V/V = 3(\lPl/ct)^2$ freezes the
sub-Planck fluctuations as the mode exits the Hubble horizon.
At last scattering ($t_{\rm LS} = 380\,\mathrm{kyr}$), the Sachs--Wolfe
plateau gives:
\begin{equation}
  \frac{\delta T}{T}\bigg|_{\rm SW}
  = \frac{1}{3}\,\frac{\delta E}{E}\left(\frac{\tPl}{t_{\rm LS}}\right)^2
  \approx 2.1\times10^{-47}.
\end{equation}

\begin{remark}
This is far below $10^{-5}$. The reason is that we have used
$\Phi \propto (t_{\rm Pl}/t_{\rm LS})^2$ directly, which applies only to
the Planck-cell mode itself (scale $\sim\lPl$). Superhorizon modes of
cosmological relevance ($k \sim 0.05\,\mathrm{Mpc}^{-1}$, corresponding to
$t_{\rm exit} \gg t_{\rm Pl}$) pick up the frozen potential
$\Phi_{\rm frozen}$ at their own horizon-crossing time, which is much larger
than $\tPl/t_{\rm LS}$. The CMB amplitude is thus set by the
\emph{comoving scale} of interest, not by the Planck-cell mode.
The correct calculation requires propagating the two-point function
$\langle T_{\mu\nu}(x)T_{\rho\sigma}(x')\rangle_{\rm ZS}$ from the
Planck epoch to last scattering using the full Boltzmann hierarchy.
\end{remark}

\subsection{Structure Scales}

The sound horizon at last scattering is:
\begin{equation}
  r_s = \int_0^{t_{\rm LS}} \frac{c_s\,dt}{a(t)} \approx 2c_s t_{\rm LS}
  = \frac{2C}{\sqrt{3}}\,t_{\rm LS} \approx 131\,\mathrm{Mpc}.
\end{equation}
The Silk damping scale (photon diffusion):
\begin{equation}
  \lambda_{\rm Silk} \approx \sqrt{r_s\,\lambda_{\rm mfp}^\gamma}
  \approx 9\,\mathrm{Mpc}.
\end{equation}
Both scales are consistent with Planck\,2018 observations.

\subsection{What Closes OP2}

\begin{openprobbox}{OP2}
Compute the two-point function
$\langle T_{\mu\nu}(x)\,T_{\rho\sigma}(x')\rangle_{\rm ZS}$
for the zero-sum constrained density operator at $t = \tPl$.
Propagate through the Boltzmann hierarchy (temperature, polarisation,
neutrino, dark matter) to last scattering.
The leading prediction of TUO is a Harrison--Zel'dovich spectrum $n_s = 1$
at $O(1)$, with a correction $\delta n_s \sim (\lPl/r_s)^2 \ll 1$
from the $q(\tPl) = -1$ curvature.
\end{openprobbox}

% ─────────────────────────────────────────────────────────────────────────────
\section{OP4: Precise Junction Time}
\label{sec:op4}
% ─────────────────────────────────────────────────────────────────────────────

\subsection{Thermalisation Condition}

The TUO--FRW junction requires both:
\begin{enumerate}
\item The quantum volume correction $\Delta V/V = 3(\lPl/ct)^2 < 1\%$,
  giving the lower bound $t_{\rm low} = \sqrt{300}\,\tPl = 17.3\,\tPl$.
\item Full QGP thermalisation, i.e.\ the optical depth $\int c/\Hmfp\,dt = 1$.
\end{enumerate}

\subsection{Two-Phase Integral}

The mean free path at $T = \TTUO$ is:
\begin{equation}
  \Hmfp(T_{\rm TUO}) = \frac{\hbar c}{\alphas(E_{\rm Pl}/2)\,k_B\TTUO}
  = 85.4\,\lPl,
\end{equation}
using $\alphas(E_{\rm Pl}/2) = 0.01055$ (1-loop QCD, PDG\,2024).

\paragraph{Phase 1} ($0 \leq t \leq t_{\rm class} = 17.3\,\tPl$): The quantum
regime. $T = \TTUO$, $\Hmfp = \Hmfp^{(0)}$.
\begin{equation}
  I_1 = \int_0^{t_{\rm class}} \frac{c}{\Hmfp^{(0)}}\,dt
  = \frac{c\,t_{\rm class}}{\Hmfp^{(0)}}
  = \frac{17.3}{85.4} = 0.2026.
\end{equation}

\paragraph{Phase 2} ($t_{\rm class} \leq t \leq \tau_{\rm th}$): Classical
FRW. The temperature cools as $T(t) = \TTUO\,(t_{\rm class}/t)^{1/2}$,
giving:
\begin{equation}
  \Hmfp(t) = \Hmfp^{(0)}\sqrt{t/t_{\rm class}}.
\end{equation}
The condition $I_1 + I_2 = 1$ becomes:
\begin{equation}
  \frac{c}{\Hmfp^{(0)}\sqrt{t_{\rm class}}}
  \int_{t_{\rm class}}^{\tau_{\rm th}} t^{-1/2}\,dt = 1 - I_1,
\end{equation}
which gives:
\begin{equation}
  \sqrt{\tau_{\rm th}} = \sqrt{t_{\rm class}} + \frac{(1-I_1)\Hmfp^{(0)}\sqrt{t_{\rm class}}}{2c}.
\end{equation}
Solving:
\begin{equation}
  \boxed{\tau_{\rm th} = 152.3\,\tPl.}
\end{equation}
Numerical verification: $\int_0^{\tau_{\rm th}} c/\Hmfp(t)\,dt = 1.000000$.

\subsection{Updated Junction Range}

\begin{table}[h]
\centering
\caption{TUO--FRW junction time bounds.}
\label{tab:junction}
\begin{tabular}{lrrl}
\toprule
Bound & Value ($\tPl$) & Method & Status \\
\midrule
Lower & 17.3 & $\Delta V/V = 1\%$ & This paper (unchanged) \\
Upper (isothermal) & 85.4 & $\Hmfp/c$ at $T_{\rm TUO}$ & Paper (2026): $\tau = 94\,\tPl$ \\
\rowcolor{lightgold}
Upper (FRW cooling) & \textbf{152.3} & FRW $T\propto t^{-1/2}$ & \textbf{This work} \\
\bottomrule
\end{tabular}
\end{table}

The FRW-cooling upper bound extends the junction range by a factor of
$152/85 \approx 1.8$ relative to the isothermal estimate. The physical
reason is that as the temperature drops, the mean free path grows, and
thermalisation takes proportionally longer.

\begin{resultbox}{OP4 -- Updated junction range}
\begin{equation}
  t_{\rm junc} \in [17.3,\ 152.3]\,\tPl.
\end{equation}
The TUO--FRW handoff does not occur at $\tPl$. Full QGP thermalisation
requires the system to expand to $\sim 150$ Planck times.
Remaining uncertainty: $\pm20\%$ from 2-loop QCD transport coefficients.
\end{resultbox}

% ─────────────────────────────────────────────────────────────────────────────
\section{The OP1--OP5 Link: A Unified Statement}
\label{sec:link}
% ─────────────────────────────────────────────────────────────────────────────

The central conceptual result of this paper is that OP1 and OP5 are not
independent. We state this as a corollary:

\begin{corollary}[OP1--OP5 equivalence]
\label{cor:op1op5}
In TUO, the following are equivalent:
\begin{enumerate}
\item The baryon asymmetry $\eta \neq 0$.
\item The leptogenesis CP asymmetry $\eps_{\rm CP} \neq 0$.
\item The CKM matrix has at least one physical CP-violating phase.
\item $\Ngen \geq 3$.
\end{enumerate}
The observed value $\eta = 6.12\times10^{-10}$ provides the lower bound
$\Ngen \geq 3$. The LEP measurement $N_\nu = 2.984 \pm 0.008$ provides the
upper bound $\Ngen \leq 3$. The intersection is $\Ngen = 3$ exactly.
\end{corollary}

This corollary does not ``derive'' $\Ngen = 3$ from the TUO axioms alone.
Rather, it shows that the TUO axioms, combined with the observational
inputs $\eta > 0$ and $N_\nu \approx 3$, uniquely determine $\Ngen = 3$.
The axioms provide the framework (zero-sum thermal bath); the observations
provide the two-sided constraint.

The remaining work to fully close both OP1 and OP5 is a single calculation:
\begin{equation}
  \text{Solve leptogenesis Boltzmann equations at } T = \TTUO,
  \text{ with } (B-L)_{\rm init} = 0 \text{ and } \Ngen = 3.
\end{equation}
The output is $\eta(M_N, Y_\nu, \delta_{\rm CP})$. Matching to $\eta_{\rm obs}$
constrains the RH neutrino sector. This is a standard---if difficult---calculation.

% ─────────────────────────────────────────────────────────────────────────────
\section{Summary of Results}
\label{sec:summary}
% ─────────────────────────────────────────────────────────────────────────────

\begin{table}[h]
\centering
\caption{Summary of new results from this analysis.}
\label{tab:summary}
\begin{tabular}{p{0.12\textwidth}p{0.35\textwidth}p{0.35\textwidth}l}
\toprule
Problem & Result & Status & Section \\
\midrule
OP5 &
$\Zcorr(g_*)$ does not select $\Ngen = 3$. Free energy strictly decreasing.
$\Zcorr$ correction $\sim 10^{-93}$/cell — negligible. &
Closed (negative result, pinch from KM + LEP) & \S\ref{sec:op5} \\[4pt]
OP1 &
$(B-L)_{\rm init} = 0$ requires leptogenesis.
$\eps_{\rm CP} \geq 1.3\times10^{-6}$, $M_N \geq 3\times10^{12}\,\mathrm{GeV}$.
KM: $\Ngen \geq 3$ required. &
Open; specific Boltzmann calculation identified & \S\ref{sec:op1} \\[4pt]
OP2 &
$\delta T/T|_{\rm cell} = 2.4\%$, suppression $\times 2400$ to CMB.
$r_s = 131\,\mathrm{Mpc}$, $\lambda_{\rm Silk} = 9\,\mathrm{Mpc}$.
Leading prediction: $n_s = 1$. &
Open; two-point function $\langle T_{\mu\nu}T_{\rho\sigma}\rangle_{\rm ZS}$ identified & \S\ref{sec:op2} \\[4pt]
OP4 &
$\tau_{\rm th} = 152.3\,\tPl$ (FRW cooling). Updated range $[17.3, 152.3]\,\tPl$.
Paper's isothermal value $94\,\tPl$ was underestimate by $\times 1.8$. &
Substantially closed; $\pm20\%$ from 2-loop QCD & \S\ref{sec:op4} \\
\bottomrule
\end{tabular}
\end{table}

% ─────────────────────────────────────────────────────────────────────────────
\section{Conclusion}
\label{sec:conc}
% ─────────────────────────────────────────────────────────────────────────────

The Z\textsubscript{corr}$(g_*)$ calculation was expected to be the key to OP5.
It is not. The thermal entropy of the pre-emergence plasma is so large---and
grows so fast with $\Ngen$---that no saddle-point correction of the form
$(\det\chi)^{-1/2}$ can create a preference for $\Ngen = 3$. This is a clean
null result, and it clarifies the structure of the problem considerably.

The resolution---a two-sided pinch from the Kobayashi--Maskawa theorem and
the LEP $Z$-boson measurement---is not an \emph{ad hoc} rescue. The KM lower
bound $\Ngen \geq 3$ is a direct consequence of the requirement that the TUO
thermal bath produce $\eta \neq 0$ via leptogenesis. In other words, TUO
plus the demand that our universe be matter-dominated forces $\Ngen \geq 3$.
The LEP upper bound then pins the value exactly.

The most important single calculation remaining across all open problems is the
full leptogenesis Boltzmann system at $T = \TTUO$ with initial condition
$(B-L) = 0$. Closing it closes both OP1 and OP5 simultaneously.

% ─────────────────────────────────────────────────────────────────────────────
\appendix
\section{Planck Units and Key Numbers}
% ─────────────────────────────────────────────────────────────────────────────

\begin{table}[h]
\centering
\caption{Key numerical results (all verified by \texttt{tuo\_open\_problems\_results.py}).}
\label{tab:numbers}
\begin{tabular}{lrl}
\toprule
Quantity & Value & Derivation \\
\midrule
$\alphas(E_{\rm Pl}/2)$ & $0.01055 \pm 0.0002$ & 1-loop QCD, PDG\,2024 \\
$\TTUO / \TPl$ & $1.1103$ & $(15/\pi^2)^{1/4}$ theorem \\
$x = \EPl/(2k_B\TTUO)$ & $0.4503$ & All modes excited ($x < 1$) \\
$n_B$ at $E = \EPl/2$ & $1.758$ & Classical bosonic regime \\
$n_F$ at $E = \EPl/2$ & $0.389$ & Partially filled fermionic \\
$\Hmfp(\TTUO)$ & $85.4\,\lPl$ & Perturbative QGP \\
$\tau_{\rm gg}$ & $1.489\,\tPl$ & Gluon-gluon scattering \\
$\tau_{\rm th}^{\rm FRW}$ & $152.3\,\tPl$ & FRW cooling integral \\
$N_\gamma$ per cell at $\TTUO$ & $0.3334$ & Stefan-Boltzmann \\
$\eps_{\rm CP}^{\rm required}$ & $\sim 1.3\times10^{-6}$ & $\eta_{\rm obs}\gstar/\kappa$ \\
$M_N^{\rm min}$ & $\sim 3\times10^{12}\,\mathrm{GeV}$ & Davidson--Ibarra \\
$\delta T/T|_{\rm cell}$ & $0.0242$ ($2.4\%$) & $1/(4\sqrt{\gstar})$ \\
Suppression to CMB & $\times 2400$ & $\sim 11$ binary e-folds \\
\bottomrule
\end{tabular}
\end{table}

% ─────────────────────────────────────────────────────────────────────────────
\section{g* Counting: Comparison of Conventions}
% ─────────────────────────────────────────────────────────────────────────────

Two conventions appear in the TUO literature:

\begin{description}
\item[Particles only (original paper):] Counts each particle species once.
  Per gen: $u(6) + d(6) + e(2) + \nu(1) = 15$ fermionic Weyl modes.
  $\gstar(3) = 28 + 3\times15\times7/8 = 67.375$, giving $E_{\rm cell} = 33.69\,\EPl$.

\item[Particles + antiparticles (this work):] Standard cosmological convention.
  Per gen: $30$ modes including antiparticles.
  $\gstar(3) = 28 + 3\times30\times7/8 = 106.75$, giving $E_{\rm cell} = 53.375\,\EPl$.
\end{description}

\emph{This paper uses the standard cosmological convention.}
The qualitative conclusion of the OP5 calculation is unchanged under either
convention: the free energy decreases monotonically, and $\Zcorr$ is
negligible.

% ─────────────────────────────────────────────────────────────────────────────
\begin{thebibliography}{9}

\bibitem{pdg2024}
  Particle Data Group (2024). \textit{Review of Particle Physics}.
  \textit{Physical Review D}, 110, 030001.

\bibitem{km1973}
  Kobayashi, M.\ and Maskawa, T. (1973).
  CP-violation in the renormalizable theory of weak interaction.
  \textit{Progress of Theoretical Physics}, 49(2), 652--657.

\bibitem{lep2006}
  ALEPH, DELPHI, L3, OPAL, SLD (2006).
  Precision electroweak measurements on the $Z$ resonance.
  \textit{Physics Reports}, 427, 257--454.

\bibitem{davidson2002}
  Davidson, S.\ and Ibarra, A. (2002).
  A lower bound on the right-handed neutrino mass from leptogenesis.
  \textit{Physics Letters B}, 535, 25--32.

\bibitem{planck2018}
  Planck Collaboration (2020).
  Planck\,2018 results VI: Cosmological parameters.
  \textit{Astronomy \& Astrophysics}, 641, A6.

\bibitem{matshaba2026}
  Matshaba, R. (2026).
  Universal Origins: The Zero-Sum Constraint --- Matrix Formulation of Cosmogenesis.
  Preprint. DOI: \href{https://doi.org/10.5281/zenodo.18827993}{10.5281/zenodo.18827993}.

\end{thebibliography}

\end{document}
